\section{Four Ways Readers Reconcile John and Luke}

The difference between John's ``believed'' and Luke's ``wondering'' has
generated several harmonizations. Each is possible; none is printed by both
Evangelists.

\textbf{Peter first wondered and later believed.} Luke describes Peter's
initial departure; John describes the other disciple's immediate response;
Peter's own Easter faith comes with the appearance tradition. This lets each
narrative subject retain his own verb.

\textbf{Both believed, but Luke names only wonder.} Wonder can accompany faith,
so Luke need not exclude an emerging conviction. Plausible psychology does not
warrant changing Luke's predicate to ``believed.''

\textbf{The other disciple believed Mary's removal report.} On this account
John 20:8 means only that the empty tomb confirmed Mary's message. Verse 9 then
explains why resurrection faith had not arisen. Against it stand the local
movement from seeing to believing and the scene's movement beyond Mary's first
inference.

\textbf{Luke 24:12 is a secondary harmonization from John.} Westcott and Hort
once judged the verse drawn from John with the other disciple omitted. The
2023 SBLGNT running text prints the verse; its edition-comparison apparatus
records the verse in Tregelles, NA28, and Robinson--Pierpont while marking it
as double-bracketed in Westcott--Hort. That control is not a manuscript
apparatus and does not settle the verse's earliest recoverable form. Even if
literary dependence were established, the printed Lukan form attributes
wonder, not belief, to Peter.

Textual history asks whether Luke 24:12 belonged to the earliest recoverable
Luke; narrative exegesis asks what the printed verse means; canonical
interpretation asks how both accounts function together. One answer cannot
silently perform all three jobs.

